% Guia do professor — Capítulo 7: Processos de Precipitação
\section{Capítulo 7 --- Processos de Precipitação}

\subsection*{Visão geral e ritmo}

Microfísica de nuvens organizada como um problema de escalas: por que a
natureza precisa de aerossol para condensar e de colisões (ou gelo) para
chover. \textbf{Ritmo sugerido: 2 aulas de 2\,h.}

\begin{itemize}
  \item \textbf{Aula 1} — escada de tamanhos, Kelvin, Köhler, ativação,
        nucleação de gelo. Casos~1--2 (15\,min; o \texttt{pyrcel} rende).
  \item \textbf{Aula 2} — difusão ($\sqrt t$ e seu gargalo),
        colisão--coalescência, Bergeron, velocidades terminais,
        Marshall--Palmer. Casos~3--4 (15\,min).
\end{itemize}

\subsection*{Objetivos de aprendizagem}

\begin{enumerate}
  \item explicar por que a atmosfera não condensa sem CCN (barreira de
        Kelvin);
  \item ler uma curva de Köhler: $r^*$, $S^*$, ativação e papel do núcleo;
  \item mostrar que difusão dá $r\propto\sqrt t$ e por que isso não faz
        chuva;
  \item comparar chuva quente (coleta) e fria (Bergeron) e saber qual domina
        onde;
  \item usar Marshall--Palmer para calcular massa, taxa de chuva e
        refletividade.
\end{enumerate}

\subsection*{Pré-requisitos a reativar}

$e_s$ sobre água/gelo (Cap.~4); tensão superficial e energia livre
(termodinâmica básica); arrasto de Stokes (mecânica dos fluidos elementar).

\subsection*{Armadilhas conceituais frequentes}

\begin{itemize}
  \item \textbf{``Nuvem cresce até virar chuva por condensação''}: o
        gargalo do $\sqrt t$ desfaz; insistir nos tempos (minutos vs.\
        dias).
  \item \textbf{Supersaturações imaginárias}: alunos assumem UR de
        110--120\,\% em nuvens; o real é $\lesssim 100{,}5$\,\% — e é por
        isso que Kelvin é barreira.
  \item \textbf{Gelo a $0\,°$C}: água super-resfriada até $-38\,°$C
        surpreende; núcleos de gelo são raros.
  \item \textbf{Marshall--Palmer como lei universal}: é climatologia de
        chuva estratiforme; convecção e neve fogem dela.
  \item \textbf{Efeito do aerossol}: mais CCN $=$ mais gotas \emph{menores}
        (não mais chuva!) — nuvem mais brilhante e chuva possivelmente
        suprimida.
\end{itemize}

\subsection*{Demonstração de baixo custo}

Garrafa PET (nuvem na garrafa) com e sem fumaça: sem núcleos quase nada
acontece; com fumaça a nuvem é instantânea — Kelvin e CCN em 2 minutos.
Vídeo de câmara de nuvens (cloud chamber) fecha a Aula 1.

\subsection*{Problemas sugeridos do livro (exercícios do Cap.~7)}

Aquecimento: UR de equilíbrio de gotas pequenas; $r^*$ e $S^*$ para núcleos
dados; velocidades terminais. Aprofundamento: tempos de crescimento por
difusão; taxa de chuva de distribuições dadas. Síntese: ``por que
semeadura de nuvens usa AgI?'' (estrutura cristalina $\approx$ gelo).
\emph{Confira a numeração na sua cópia (v1.02b).}

\subsection*{Uso do notebook \texttt{cap07\_precipitacao.ipynb}}

\begin{center}
\begin{tabular}{lll}
\hline
Caso & Conteúdo & Momento sugerido \\
\hline
1 & Curvas de Köhler; $r^*(r_d)$, $S^*(r_d)$ & Aula 1, após Köhler \\
2 & \texttt{pyrcel}: ativação vs.\ updraft & fim da Aula 1 \\
3 & Difusão × coleta: o gargalo e a avalanche & Aula 2, após crescimento \\
4 & Marshall--Palmer: momentos, $R$, $Z$, relação $Z$--$R$ & fim da Aula 2 \\
\hline
\end{tabular}
\end{center}

O Caso~2 usa um pacote de pesquisa de verdade (\texttt{pyrcel}) — vale
contar aos alunos que é o mesmo tipo de módulo que roda dentro de modelos
climáticos. O Caso~4 fecha preparando o radar do Cap.~8.

\subsection*{Exercícios complementares e soluções}

Nas notas: T1--T5 (Kelvin por energia livre, extremos de Köhler, tempos de
difusão, Stokes, momentos de M--P) e N1--N4. Soluções em
\texttt{cap07\_solucoes.ipynb} (\emph{não distribuir}): T2 e T5 com
\texttt{sympy} (extremos e integrais $\Gamma$), N4 verificando $Z = aR^b$.
Pares: T2$\leftrightarrow$N1, T5$\leftrightarrow$N4.

\subsection*{Conexões com o resto do curso}

$e_s$ água/gelo $\to$ Cap.~4; correntes ascendentes que sustentam gotas
$\to$ Caps.~5 e 14; $Z$--$R$ $\to$ radar (Cap.~8); aerossol--nuvem $\to$
clima (Cap.~21); granizo $\to$ Cap.~15.
